Modern applications consist of several interconnected services deployed on
shared infrastructures such as public clouds, where physical resources such as
processor and memory are shared between multiple applications to optimize
resource usage and reduce costs. With this paradigm, strong isolation techniques
are needed to preserve the confidentiality and integrity of data and
computations. \acp{TEE} tackle this problem by relying on hardware extensions to
the underlying computing architecture to provide strong isolation and access
control both in the main processor and in memory. Code and data running inside a
\ac{TEE} are then shielded from the rest of the system, and even privileged
software such as the hypervisor is not allowed to access the secure environment,
thereby significantly reducing the overall attack surface. In recent years, the
community has adopted the term \emph{confidential computing} to define \ac{TEE}
workloads.

In this Ph.D.\ dissertation, we analyze \acp{TEE} from a practical standpoint,
focusing on how to securely and effectively integrate them into novel or
existing applications and protocols. We study two use cases, namely smart homes
and \ac{V2X} credential management, where \acp{TEE} can be utilized to enhance
the security posture of the system while providing additional benefits compared
to traditional approaches, such as increased efficiency and scalability.
Furthermore, we raise the attention to non-obvious practical challenges of this
technology: We investigate the interface between an application running in a
\ac{TEE} and the untrusted world, showing how attackers can manipulate the
notion of time perceived inside the \ac{TEE} to compromise the integrity of the
application. Besides, we analyze the trust relationships between customers who
want to leverage \acp{TEE} in the cloud and their providers, highlighting some
inconsistencies with the \ac{TEE} threat model. We conclude that this technology
can become a cornerstone of modern applications, but its integration requires
careful design choices from developers and administrators, as well as proper
support from the infrastructure providers.